\chapter{Interfacing with External Systems}
\label{ch:external}

A typical Kiss application is more than a self-contained box: it calls third-party APIs, accepts inbound webhooks, and increasingly both exposes tools to AI assistants and consumes tools from them.  This chapter covers the seven classes the framework provides for that work:

\begin{itemize}
  \item \texttt{RestClient}\index{RestClient} --- make HTTP calls \emph{out} to external services.
  \item \texttt{RestServerBase}\index{RestServerBase} --- accept HTTP calls \emph{in} from external services that don't speak Kiss's JSON-RPC.
  \item \texttt{RestServerExample}\index{RestServerExample} --- a working template that extends \texttt{RestServerBase}.
  \item \texttt{MCPServerBase}\index{MCPServerBase} --- expose the application as a Model Context Protocol (MCP) server, the standard way for AI assistants to discover and invoke tools.
  \item \texttt{MCPServerExample}\index{MCPServerExample} --- a working template that extends \texttt{MCPServerBase}.
  \item \texttt{MCPClientBase}\index{MCPClientBase} --- consume tools from a \emph{remote} MCP server, the mirror image of \texttt{MCPServerBase}.
  \item \texttt{MCPClientExample}\index{MCPClientExample} --- a working template that extends \texttt{MCPClientBase}.
\end{itemize}

\noindent
The three \texttt{*Base}/\texttt{*Example} pairs are intentionally parallel: \texttt{*Base} carries the protocol boilerplate and lives in \texttt{src/main/core/}; \texttt{*Example} is a copyable template under \texttt{src/main/precompiled/} that you adapt for each endpoint or remote server you work with.

\noindent
The \texttt{RestClient} side is symmetric --- it is the same set of operations any HTTP client offers, just bound to Kiss's JSON utilities.  The server side has more nuance, because there is already a built-in JSON-RPC entry point (\texttt{/rest}) used by the Kiss frontend; the classes here are for the cases where that is not the protocol the caller is using.

\begin{notebox}
Use a normal Kiss JSON-RPC service (Chapter~\ref{ch:services}) for everything called from the Kiss frontend.  Reach for the classes in this chapter only when an \emph{external}, non-Kiss client is the caller --- typically because that client demands a specific URL, HTTP method, header set, authentication scheme, or response format that does not fit Kiss's JSON-RPC convention.
\end{notebox}

\section{RestClient --- Calling External APIs}
\label{sec:ext-restclient}

\texttt{RestClient} is an HTTP client for making REST calls to external APIs.  It supports JSON, string, XML, and streaming response types.

\begin{lstlisting}[style=groovy]
import org.kissweb.RestClient

// JSON GET request
JSONObject result = RestClient.jsonCall(
    "GET", "https://api.example.com/users/123")

// JSON POST with body
JSONObject body = new JSONObject()
body.put("name", "Alice")
body.put("email", "alice@example.com")
JSONObject result = RestClient.jsonCall(
    "POST", "https://api.example.com/users", body)

// String response
String html = RestClient.strCall(
    "GET", "https://example.com/page")

// XML response
Document xml = RestClient.xmlCall(
    "GET", "https://api.example.com/data.xml")

// Streaming response
InputStream stream = RestClient.streamCall(
    "GET", "https://example.com/large-file.zip")
\end{lstlisting}

\noindent
Inside a Kiss service the same pattern applies.  The HTTP call is just a piece of the service's normal logic:

\begin{lstlisting}[style=groovy]
import org.kissweb.RestClient

void fetchWeather(JSONObject injson, JSONObject outjson,
                  Connection db, ProcessServlet servlet) {
    String city = injson.getString("city")
    RestClient rc = new RestClient()
    JSONObject result = rc.jsonGet(
        "https://api.example.com/weather?city=" + city)
    outjson.put("temperature", result.getDouble("temp"))
}
\end{lstlisting}

\subsection{Advanced Configuration}

\begin{lstlisting}[style=groovy]
RestClient client = new RestClient()

// Proxy support
client.setProxy("proxy.corp.com", 8080)

// TLS client certificate
client.setTLSKey("keystore.p12", "password")

// Retry policy for transient failures
client.setRetryPolicy(3, 1000)  // 3 retries, 1s delay
\end{lstlisting}

\begin{longtable}{P{4.8cm}P{5.7cm}}
\toprule
\textbf{Method} & \textbf{Description} \\
\midrule
\endhead
\texttt{jsonCall(method, url)} & HTTP call, parse JSON response. \\
\texttt{jsonCall(method, url, body)} & HTTP call with JSON body. \\
\texttt{strCall(method, url)} & HTTP call, return string. \\
\texttt{xmlCall(method, url)} & HTTP call, parse XML response. \\
\texttt{streamCall(method, url)} & HTTP call, return InputStream. \\
\texttt{setProxy(host, port)} & Configure HTTP proxy. \\
\texttt{setTLSKey(path, pw)} & Set client certificate. \\
\texttt{setRetryPolicy(n, delay)} & Set retry count and delay. \\
\bottomrule
\end{longtable}

\section{RestServerBase --- Receiving External REST Calls}
\label{sec:ext-restserverbase}

When an \emph{external} system needs to call into your application (a webhook from a payment processor, an inbound notification from a third-party API, a status callback from a remote job, etc.), that system will not speak Kiss's JSON-RPC protocol.  For these cases the framework provides a thin base class, \texttt{org.kissweb.RestServerBase}, that lets you expose a conventional servlet-style REST endpoint.

\subsection{What \texttt{RestServerBase} Provides}

\texttt{RestServerBase} extends \texttt{HttpServlet} and adds a small set of static helpers so the endpoint code stays focused on application logic:

\begin{itemize}
  \item \texttt{basicAuthenticate(request, response, realm, user, password)} --- validates an HTTP Basic \texttt{Authorization} header against an expected username and password, and sets a 401 response on failure.
  \item \texttt{getBodyString(request)} --- read the full request body as a \texttt{String}.
  \item \texttt{getBodyJson(request)} --- read the body as a \texttt{JSONObject} (returns \texttt{null} if the content type is not \texttt{application/json}).
  \item \texttt{getBodyBinary(request)} --- read the body as a \texttt{byte[]}.
  \item \texttt{setResponse(response, jsonObject)} --- write a JSON response with the correct content type and UTF-8 encoding.
  \item \texttt{setResponse(response, contentType, string)} and \texttt{setResponse(response, contentType, byte[])} --- write text or binary responses.
  \item \texttt{isTextContentType(contentType)} / \texttt{isBinaryContentType(contentType)} --- best-guess classification of an HTTP content type.
\end{itemize}

Override any of \texttt{doGet}, \texttt{doPost}, \texttt{doPut}, or \texttt{doDelete} to handle the corresponding HTTP method; everything else follows standard Jakarta Servlet conventions.

\subsection{Configuration --- No \texttt{web.xml} Required}

The \texttt{@WebServlet} annotation is the \emph{only} configuration required to register the endpoint.  The URL the annotation declares (e.g.\ \texttt{/myservice}) is appended to the application's context path; no \texttt{<servlet>} or \texttt{<servlet-mapping>} entry in \texttt{web.xml} is needed.  Servlet 6.x annotation scanning discovers the class at startup automatically.

\begin{notebox}
The framework's \texttt{src/main/core/WEB-INF/web.xml} only declares the framework's startup listener and a global CORS filter mapped to \texttt{/*}.  Your endpoint inherits CORS handling for free; you do not add anything to \texttt{web.xml} when adding new external endpoints.
\end{notebox}

\subsection{Why It Must Live Under \texttt{precompiled/}}

Servlet annotation scanning happens once, when the servlet container starts.  Classes that the container needs to discover at startup must therefore be compiled into the WAR ahead of time --- they cannot be supplied by Kiss's runtime compiler, which only loads files under \texttt{src/main/backend/} on demand.  For this reason, every class extending \texttt{RestServerBase} must be placed under \texttt{src/main/precom\-piled/}, where \texttt{./bld -v build} compiles it into \texttt{work/exploded/WEB-INF/classes/} before Tomcat starts.

The practical consequence: after adding or modifying a \texttt{RestServerBase} subclass, run \texttt{./bld -v build} and restart the server.  The hot-reload behavior that applies to backend services does \emph{not} apply here.

\begin{tipbox}
\texttt{@MultipartConfig} is only needed when the endpoint accepts \texttt{multi\-part/form-data} uploads.  Omit it for plain JSON or text endpoints.
\end{tipbox}

\section{RestServerExample --- The Template}
\label{sec:ext-restserverexample}

A working template lives at \texttt{src/main/precom\-piled/com/mycompany/Rest\-Server\-Example.java}.  Copy it, rename it, and edit it for each external endpoint you need to expose.

\begin{lstlisting}[style=java]
package com.mycompany;

import jakarta.servlet.annotation.WebServlet;
import jakarta.servlet.annotation.MultipartConfig;
import jakarta.servlet.http.HttpServletRequest;
import jakarta.servlet.http.HttpServletResponse;
import org.kissweb.RestServerBase;
import org.kissweb.json.JSONObject;
import java.io.IOException;

@WebServlet(urlPatterns="/myservice")
@MultipartConfig
public class RestServerExample extends RestServerBase {

    @Override
    protected void doPost(HttpServletRequest request,
            HttpServletResponse response) throws IOException {
        if (!basicAuthenticate(request, response, null,
                               "username", "password"))
            return;

        JSONObject in = getBodyJson(request);
        // ... process the request ...
        JSONObject out = new JSONObject();
        out.put("ok", true);
        setResponse(response, out);
    }
}
\end{lstlisting}

\noindent
The shipped example also stubs out \texttt{doGet}, \texttt{doPut}, and \texttt{doDelete} as override points; delete the methods you don't need.

\section{MCPServerBase --- Exposing Tools to AI Assistants}
\label{sec:ext-mcpserverbase}

The Model Context Protocol (MCP) is the open protocol AI assistants such as Claude Desktop and Claude Code use to talk to external tools and data.  An MCP server advertises a list of tools (each with a JSON-Schema for its arguments); the assistant chooses when to call them and feeds the results back into the model's reasoning.

\texttt{MCPServerBase} is the framework's MCP server base class.  It extends \texttt{RestServer\-Base}, so it inherits the same body-reading and response-writing helpers, the same \texttt{@WebServlet}-based registration model, and the same precompiled-directory requirement.  On top of that it implements all of the wire-protocol boilerplate --- the JSON-RPC 2.0 envelope, the \texttt{initialize} handshake, \texttt{ping}, notifications, and error replies --- so a subclass only has to declare its tools and run them.

\subsection{The Wire Protocol}

Every MCP exchange is a JSON-RPC 2.0 message:

\begin{lstlisting}[style=java]
// request
{ "jsonrpc": "2.0", "id": 1, "method": "tools/list",
    "params": {} }

// response
{ "jsonrpc": "2.0", "id": 1, "result": { "tools": [ ... ] } }

// notification (no id, no response expected)
{ "jsonrpc": "2.0", "method": "notifications/initialized" }
\end{lstlisting}

\noindent
\texttt{MCPServerBase.doPost} parses the envelope, dispatches on the \texttt{method} field, and either returns a JSON response (for requests) or a 202 Accepted with no body (for notifications).  It is declared \texttt{final} --- subclasses extend behavior through hooks rather than by overriding the dispatcher.

\subsection{The Three Methods That Matter}

A minimum-viable MCP server implements three methods, all of which \texttt{MCPServerBase} dispatches automatically:

\begin{itemize}
  \item \texttt{initialize} --- handshake.  Client sends its protocol version and info; server responds with its protocol version, server info, and \emph{capabilities} (which feature areas it supports).
  \item \texttt{tools/list} --- catalog.  Server returns an array of tools, each with \texttt{name}, \texttt{description}, and an \texttt{inputSchema} (a JSON Schema describing the expected arguments).  The model uses these to decide when and how to call each tool.
  \item \texttt{tools/call} --- invocation.  Server runs the named tool with the supplied arguments and returns a list of \emph{content blocks} (most commonly a single \texttt{type:"text"} block).
\end{itemize}

\noindent
\texttt{MCPServerBase} also handles \texttt{ping} (the MCP keep-alive) and the standard \texttt{notifications/*} client-to-server notifications without any subclass involvement.

\subsection{Subclass Hooks}

A subclass implements four required hooks and may override three optional ones:

\begin{longtable}{P{4.8cm}P{5.7cm}}
\toprule
\textbf{Method} & \textbf{Purpose} \\
\midrule
\endhead
\multicolumn{2}{l}{\emph{Required}} \\
\texttt{getServerName()} & Server name reported during \texttt{initialize}. \\
\texttt{getServerVersion()} & Server version reported during \texttt{initialize}. \\
\texttt{listTools()} & Return the tool catalog as a \texttt{JSONArray}. \\
\texttt{callTool(name, args)} & Execute a tool, return its result object. \\
\midrule
\multicolumn{2}{l}{\emph{Optional}} \\
\texttt{authenticate(req, resp)} & Gate every request.  Default: allow all. \\
\texttt{getCapabilities()} & Capabilities advertised at \texttt{initialize}.  Default: tools only. \\
\texttt{dispatchExtension(\allowbreak method, params)} & Handle methods beyond the built-in set (\texttt{resources/list}, \texttt{prompts/get}, etc.). Default: return null, which produces a method-not-found error. \\
\bottomrule
\end{longtable}

\subsection{Helpers}

\texttt{MCPServerBase} exposes static helpers for the common shapes of MCP data:

\begin{itemize}
  \item \texttt{textBlock(text)} --- builds a \texttt{\{"type":"text","text":...\}} content block.
  \item \texttt{toolResult(text)} --- builds a complete \texttt{tools/call} result containing one text block.
  \item \texttt{toolError(message)} --- builds a \texttt{tools/call} result with \texttt{isError: true}.  Use this for tool-specific failures the model should see and react to (see below).
  \item \texttt{buildSchema(title, required, names, types, descriptions)} --- builds a JSON-Schema object for a tool's \texttt{inputSchema} from parallel arrays.
\end{itemize}

\noindent
The class also exposes JSON-RPC error code constants (\texttt{ERR\_PARSE}, \texttt{ERR\_INVALID\_REQUEST}, \texttt{ERR\_METHOD\_NOT\_FOUND}, \texttt{ERR\_INVALID\_PARAMS}, \texttt{ERR\_INTERNAL}) and a protected \texttt{sendError} method for extension handlers that need to fail fast.

\subsection{Tool Errors vs.\ Protocol Errors}

There are two distinct error channels and they mean different things:

\begin{itemize}
  \item \emph{Protocol errors} --- malformed JSON, unknown method, missing parameters --- are returned as JSON-RPC errors (\texttt{-32700} parse, \texttt{-32600} invalid request, \texttt{-32601} method not found, \texttt{-32603} internal).  The client treats these as transport failures.
  \item \emph{Tool errors} --- ``the database lookup failed,'' ``the file you asked about doesn't exist'' --- are returned as a normal \texttt{tools/call} \emph{result} with \texttt{isError: true}.  The model \emph{sees} the failure and can react to it; it does not see protocol errors.
\end{itemize}

\noindent
Use \texttt{toolError} (the latter channel) for anything you want the AI to be able to recover from, retry, or report back to the user.  Anything thrown out of \texttt{callTool} becomes a JSON-RPC \texttt{ERR\_INTERNAL} response, which the model does not see.

\subsection{What MCPServerBase Does Not Implement}

To stay focused, the base class omits several MCP features.  All of them layer on top of the same JSON-RPC envelope and can be added by overriding \texttt{dispatchExtension} and \texttt{getCapabilities}:

\begin{itemize}
  \item \emph{Resources} --- read-only data the model can request (files, database rows, etc.).  Implement \texttt{resources/list} and \texttt{resources/read}; advertise the \texttt{resources} capability.
  \item \emph{Prompts} --- reusable prompt templates the host UI can offer to the user.  Implement \texttt{prompts/list} and \texttt{prompts/get}; advertise the \texttt{prompts} capability.
  \item \emph{Sampling} --- server-initiated requests asking the client's model to generate something on the server's behalf.
  \item \emph{SSE streaming} --- responses streamed back as Server-Sent Events instead of a single JSON body.  Useful for long-running tools that produce incremental output.
  \item \emph{Sessions} --- the \texttt{Mcp-Session-Id} header lets a server keep per-client state across requests.
\end{itemize}

\section{MCPServerExample --- The Template}
\label{sec:ext-mcpserverexample}

A working template lives at \texttt{src/main/precom\-piled/com/mycompany/MCPServer\-Example.java}.  It extends \texttt{MCPServerBase}, exposes an endpoint at \texttt{/mcp}, and ships with two example tools (\texttt{echo} and \texttt{add}).  Because all of the protocol work is in the base class, the entire example is just the four hooks:

\begin{lstlisting}[style=java]
@WebServlet(urlPatterns="/mcp")
public class MCPServerExample extends MCPServerBase {

    @Override protected String getServerName() {
        return "kiss-mcp-example";
    }
    @Override protected String getServerVersion() {
        return "1.0.0";
    }

    @Override
    protected JSONArray listTools() {
        JSONArray tools = new JSONArray();

        JSONObject add = new JSONObject();
        add.put("name", "add");
        add.put("description",
            "Add two numbers and return the sum.");
        add.put("inputSchema", buildSchema(
                "Add input",
                // required
                new String[]{"a", "b"},
                new String[]{"a", "b"},
                new String[]{"number", "number"},
                new String[]{"First addend",
                    "Second addend"}));
        tools.put(add);
        // ... echo tool elided ...
        return tools;
    }

    @Override
    protected JSONObject callTool(String name,
        JSONObject args) {
        switch (name) {
            case "add": {
                double a = args.getDouble("a", 0.0);
                double b = args.getDouble("b", 0.0);
                return toolResult(String.valueOf(a + b));
            }
            // ... echo case elided ...
            default:
                return toolError("Unknown tool: " + name);
        }
    }
}
\end{lstlisting}

\noindent
The two methods stay in sync by convention: every entry in \texttt{listTools} needs a matching \texttt{case} in \texttt{callTool}, and vice versa.

\subsection{Build and Authentication}

Like \texttt{RestServerExample}, \texttt{MCPServerExample} relies on annotation scanning, lives under \texttt{src/main/pre\-compiled/}, and requires a build + restart to pick up changes:

\begin{lstlisting}[style=shell]
./bld -v build
\end{lstlisting}

\noindent
The shipped example does not authenticate.  Most MCP clients support adding an \texttt{Authorization} header to outgoing requests, so basic auth drops in by overriding \texttt{authenticate}:

\begin{lstlisting}[style=java]
@Override
protected boolean authenticate(HttpServletRequest req,
                               HttpServletResponse resp) {
    return basicAuthenticate(req, resp, null, "user",
        "secret");
}
\end{lstlisting}

\noindent
For OAuth 2.1 bearer-token authentication --- which the MCP specification mandates --- no override is needed.  Set \texttt{OAuthAuthorizationServer} in \texttt{application.ini} and every \texttt{MCPServerBase} subclass automatically validates incoming tokens against the configured authorization server.  Chapter~\ref{ch:oauth} covers this in detail.

\section{MCPClientBase --- Consuming a Remote MCP Server}
\label{sec:ext-mcpclientbase}

Where \texttt{MCPServerBase} \emph{exposes} tools to an AI assistant, \texttt{MCPClientBase} sits on the other end of the same wire: it \emph{consumes} the tools a remote MCP server offers.  It is the mirror image of the server base class --- the two speak the identical JSON-RPC 2.0 protocol (protocol version \texttt{2025-06-18}) and interoperate directly, so one Kiss app can call another Kiss app's \texttt{MCPServerBase} endpoint with nothing between them but HTTP.

Two differences from the server side are worth noting up front:

\begin{itemize}
  \item It is \emph{not} a servlet.  There is no \texttt{@WebServlet} annotation and no inbound URL; you \emph{instantiate or subclass} it from a backend service and call its methods.  Because it does not depend on annotation scanning, \texttt{MCPClientBase} itself lives in \texttt{src/main/core/} as ordinary framework code rather than under \texttt{precompiled/}.  (The shipped \emph{template} still lives under \texttt{precompiled/} so it sits beside the other examples.)
  \item HTTP is done through \texttt{RestClient} (Section~\ref{sec:ext-restclient}), so timeouts, proxying, and TLS behave as they do for any other outbound call.
\end{itemize}

\subsection{The Lazy Handshake}

MCP requires a client to perform an \texttt{initialize} handshake --- and send a \texttt{notifications/initialized} follow-up --- before any other request.  \texttt{MCPClientBase} runs this automatically the first time you call any operation, so the common case needs no ceremony:

\begin{lstlisting}[style=groovy]
import org.kissweb.MCPClientBase
import org.kissweb.restServer.MainServlet

class MyMcpClient extends MCPClientBase {
    protected String getServerUrl() {
        return (String) MainServlet.getEnvironment(
            "RemoteMcpUrl")
    }
}

MyMcpClient client = new MyMcpClient()
// runs initialize on first use
JSONArray tools = client.listTools()
String reply = client.callToolText("echo",
        new JSONObject().put("message", "hi"))
\end{lstlisting}

\noindent
Call \texttt{initialize()} explicitly only when you want the server's advertised capabilities --- \texttt{getServerInfo()}, \texttt{getServerCapabilities()}, \texttt{getNegotiatedProtocolVersion()} --- before making any tool calls.

\begin{tipbox}
\texttt{MainServlet.getEnvironment()} returns \texttt{Object}, not \texttt{String}.  Groovy assigns it to a \texttt{String} variable without complaint, but a \texttt{Java} subclass of \texttt{MCPClientBase} must cast it: \texttt{return (String) MainServlet.getEnvironment("RemoteMcpUrl");}
\end{tipbox}

\subsection{Subclass Hooks}

A subclass implements one required hook and may override four optional ones:

\begin{longtable}{P{4.8cm}P{5.7cm}}
\toprule
\textbf{Method} & \textbf{Purpose} \\
\midrule
\endhead
\multicolumn{2}{l}{\emph{Required}} \\
\texttt{getServerUrl()} & Full URL of the remote MCP endpoint.  Per Kiss convention, read from \texttt{application.ini}. \\
\midrule
\multicolumn{2}{l}{\emph{Optional}} \\
\texttt{getClientName()} & Client name reported during \texttt{initialize}.  Default provided. \\
\texttt{getClientVersion()} & Client version reported during \texttt{initialize}.  Default provided. \\
\texttt{getOAuthProviderName()} & Name of an \texttt{[OAuthClient \emph{name}]} provider to obtain bearer tokens from.  Default: none. \\
\texttt{applyAuthentication(\allowbreak headers)} & Attach authentication headers to each request.  Default: an OAuth bearer token when a provider name is set, otherwise nothing. \\
\texttt{getClientCapabilities()} & Client capabilities sent at \texttt{initialize}.  Default: none. \\
\bottomrule
\end{longtable}

\subsection{Operations}

\begin{longtable}{P{4.8cm}P{5.7cm}}
\toprule
\textbf{Method} & \textbf{Description} \\
\midrule
\endhead
\texttt{initialize()} & Run the handshake (once); return the server's \texttt{initialize} result. \\
\texttt{listTools()} & Fetch the tool catalog as a \texttt{JSONArray}; \texttt{nextCursor} pagination is followed automatically. \\
\texttt{callTool(name, args)} & Call a tool; return the raw result object. \\
\texttt{callToolText(name, args)} & Call a tool and return its text content; throws on a tool error. \\
\texttt{ping()} & MCP keep-alive. \\
\texttt{call(method, params)} & Arbitrary JSON-RPC request, for methods without a wrapper. \\
\texttt{notify(method, params)} & Arbitrary JSON-RPC notification. \\
\texttt{setTimeouts(connect, request)} & Connect and request timeouts (defaults 30\,s / 60\,s). \\
\bottomrule
\end{longtable}

\subsection{Reading Tool Results}

A \texttt{tools/call} result is a list of content blocks plus an optional \texttt{isError} flag.  Two static helpers cover the common cases:

\begin{itemize}
  \item \texttt{isError(result)} --- \texttt{true} if the tool reported \texttt{isError: true}.
  \item \texttt{textOf(result)} --- the \texttt{text} of every \texttt{type:"text"} block, joined with newlines (non-text blocks are skipped).
\end{itemize}

\noindent
\texttt{callToolText} combines them: it returns \texttt{textOf(result)} but throws \texttt{MCPClientException} when \texttt{isError} is set --- the convenient ``give me the text, fail on error'' path.  Use \texttt{callTool} plus the helpers when you need the raw blocks or want to handle a tool error without an exception.

\subsection{Tool Errors vs.\ Protocol Errors}

This is the client-side mirror of the two error channels described for the server (Section~\ref{sec:ext-mcpserverbase}).  A \emph{protocol} failure --- a JSON-RPC error from the server, a non-2xx HTTP status, an unparseable body, or a transport failure --- is raised as \texttt{MCPClientBase.MCPClientException} (a \texttt{Runtime\-Exception}); when the failure carries a JSON-RPC error code it is available via \texttt{getCode()} (otherwise \texttt{0}).  A \emph{tool} error the model is meant to see arrives instead as a normal result with \texttt{isError: true}, which \texttt{callTool} returns without throwing.

\subsection{Authentication}

Authentication is the symmetric counterpart of the server's.  The server \emph{validates} incoming bearer tokens (the OAuth 2.1 resource-server role); the client \emph{obtains} them (the relying-party role, Chapter~\ref{ch:oauth-client}).  Point the client at a configured \texttt{[OAuthClient \emph{name}]} provider by overriding \texttt{getOAuthProviderName()}, and every request carries a transparently-refreshed \texttt{Authorization: Bearer} header --- discovery, dynamic registration, and token refresh are all handled by the OAuth client.

When no usable token is held and an interactive login is required, \texttt{OAuthAuthorizationRequiredException} propagates out of the call.  That is a normal control-flow signal, not a failure: catch it where you make the call and start the browser redirect flow.

\begin{lstlisting}[style=groovy,basicstyle=\ttfamily\scriptsize]
import org.kissweb.oauth.client.OAuthAuthorizationRequiredException
import org.kissweb.oauth.client.OAuthClient

try {
    reply = client.callToolText("echo", args)
} catch (OAuthAuthorizationRequiredException e) {
    String url = OAuthClient.forProvider("myprovider")
            .beginAuthorization("http://localhost:8001")
    // redirect the browser to 'url'; retry after the
    // callback completes.
}
\end{lstlisting}

\noindent
For any other scheme --- HTTP basic auth, a static API token, a custom header --- override \texttt{applyAuthentication(headers)} directly instead:

\begin{lstlisting}[style=java]
@Override
protected void applyAuthentication(JSONObject headers) {
    headers.put("Authorization", "Bearer "
            + (String) MainServlet.getEnvironment(
                "RemoteMcpToken"));
}
\end{lstlisting}

\subsection{What MCPClientBase Does Not Implement}

Mirroring the server base class, the client omits the same set of features, each of which can be added incrementally on top of the same JSON-RPC envelope: resources, prompts, sampling, server-initiated requests, SSE (\texttt{text/event-stream}) streaming responses, and \texttt{Mcp-Session-Id} session management.  (Because \texttt{RestClient} surfaces only the response body, not its headers, session management in particular would require reading response headers.)

\section{MCPClientExample --- The Template}
\label{sec:ext-mcpclientexample}

A working template lives at \texttt{src/main/precom\-piled/com/mycompany/MCPClient\-Example.java}.  It extends \texttt{MCPClientBase}, reads the remote URL from \texttt{application.ini}, and provides a static \texttt{demo()} helper that lists the remote server's tools and calls its \texttt{echo} tool --- the natural counterpart to the \texttt{echo} tool in \texttt{MCPServerExample}:

\begin{lstlisting}[style=java]
package com.mycompany;

import org.kissweb.MCPClientBase;
import org.kissweb.json.JSONArray;
import org.kissweb.json.JSONObject;
import org.kissweb.restServer.MainServlet;

public class MCPClientExample extends MCPClientBase {

    @Override
    protected String getServerUrl() {
        return (String) MainServlet.getEnvironment(
            "RemoteMcpUrl");
    }

    public static String demo() {
        MCPClientExample client = new MCPClientExample();

        // runs initialize
        JSONArray tools = client.listTools();
        JSONObject args = new JSONObject();
        args.put("message", "Hello from the Kiss MCP client");
        return client.callToolText("echo", args);
    }
}
\end{lstlisting}

\noindent
The template lives under \texttt{precompiled/} so it ships beside \texttt{MCP\-Server\-Example}; like any class there it is compiled by \texttt{./bld -v build} and requires a restart to pick up changes.  The file also carries commented-out blocks showing how to wire OAuth authentication (via \texttt{getOAuthProviderName()}) or a custom header (via \texttt{applyAuthentication()}), as described above.  Call \texttt{demo()} from any backend service:

\begin{lstlisting}[style=groovy]
String reply = com.mycompany.MCPClientExample.demo()
\end{lstlisting}
